% RESEARCH_DRAFT_SHA256: fc0b127081af2db441e3514698c4c1c571b6d042268071323aab0e1805ab0bf0
% Layout and source-reference wrappers expanded for the web renderer.
\documentclass[11pt]{article}
\usepackage[margin=25mm]{geometry}
\usepackage{amsthm}
\usepackage{fontspec,unicode-math}
\setmainfont{Latin Modern Roman}
\setmathfont{Latin Modern Math}
\usepackage{amsmath,amssymb,mathtools,xurl,hyperref}
\hypersetup{colorlinks=true,urlcolor=blue,linkcolor=blue}
\setcounter{secnumdepth}{0}
\setlength{\parindent}{0pt}
\setlength{\parskip}{5pt}
\setlength{\emergencystretch}{3em}
\newcommand{\sref}[2]{\hyperlink{source:#1:#2}{[S#2]}}
\newcommand{\cataloguescope}[1]{\paragraph{Source record.}#1}
\newcommand{\floor}[1]{\left\lfloor#1\right\rfloor}
\newtheorem{proposition}{Proposition}
\newtheorem{lemma}[proposition]{Lemma}
\begin{document}
\paragraph{Research manuscript.} OpenAI \texttt{gpt-6-astra}, reasoning \texttt{ultra}; 11 October 2026.
\paragraph{Preserved source body.} The complete catalogue mathematical text and sources follow with layout and source links expanded, including source qualifications and its historical preparation notes. New research begins at the section headed Research attempt.

\section{Sensor-Actuator Co-Design and Distributed Controllability Algorithms}\label{statement}

\subsection{Definitions and mathematical statement}

\paragraph{Shared notation.}Standard number systems are $\mathbb N,\mathbb Z,\mathbb Q,\mathbb R,\mathbb C$. Specialized notation and index ranges follow the cited source and the definitions below.


A controllability matrix for $\dot x=Ax+Bu$ is $[B,AB,\ldots,A^{n-1}B]$; full row rank is the linear-system controllability criterion. A controllability Gramian measures reachable-state energy. ADMM alternates minimization and dual updates for a split constrained optimization problem.
\paragraph{Record scope and qualifications.}
The catalogue title is retained as metadata. Named certificates, contracts, bounds or frameworks in that title are not asserted as results of the original source extract.
The budget, cost, system and PDE classes, local-information model, and ADMM formulation are not given. The entire bundle is preserved, without claiming one subquestion covers the others.

Source-record limitation: The sensor/control cost models, distributed task, dynamics, and PDE-ADMM formulation are absent; the bundle cannot be recovered as a formal target.
\paragraph{Statement recorded in the supplied source.}
Given a sensing or actuation budget and a control-cost functional, how should sensors or actuators be selected to minimize cost? Is there a local distributed algorithm for the rank or controllability-volume problem? Characterize when adding or removing links from the system dynamics or communication network improves sensing or controllability rather than amplifying disturbances. Prove convergence of ADMM to the optimal solution in the PDE setting. \href{https://doi.org/10.4171/owr/2024/56}{[S2]}
\paragraph{Source status.}
The supplied catalogue labels this source record open; that metadata was checked against the supplied snapshot on 2026-10-11. The label is preserved for provenance and is not an independent verification of present open status or mathematical correctness. The catalogue literature assessment is dated 2026-08-22.
\subsection{Short English statement}
Optimize sensing and actuation, understand link changes, find distributed controllability algorithms, and prove the stated PDE-optimization convergence.
\subsection{Sources}
\begin{itemize}
\item[S1] Ulam AI and the UnsolvedMath Contributors, UnsolvedMath: A Curated Collection of Open Mathematics Problems, catalogue entry. Catalogue curation and metadata: CC BY 4.0. The source extract is retained; definitions, English summaries and qualifications are editorial additions.
\url{https://huggingface.co/datasets/ulamai/UnsolvedMath}.
\item[S2] Mini-Workshop: High-Dimensional Control Problems and Mean-Field Equations with Applications in Machine Learning (2025), Oberwolfach Reports, DOI 10.4171/owr/2024/56, pp. 3211--3254
\url{https://doi.org/10.4171/owr/2024/56}.
\item Source-specific Oberwolfach Reports licences, including CC BY-SA 4.0 where applicable, are retained; see THIRD\_PARTY\_NOTICES.md.
\end{itemize}
\paragraph{Preparation and review.}
The original source and UnsolvedMath contributors retain attribution for the source extract. Codex (GPT-6; exact serving revision not exposed) assisted with definitions, English summaries, scope notes and formatting. No new human mathematical authorship or independently verified solution is asserted. Catalogue identity checked against the complete supplied registry and current repository. Distinct upstream records, including repeated questions and broad or damaged extracts, are retained for complete coverage of the source CS category. A repeated record is not a new independent conjecture. Complete mathematical-equivalence and current-literature review has not been performed. Preparation and final review follow contribution prompts C and D, with source gaps disclosed. Structural and rendering checks do not establish a new theorem or open status.
\label{end:statement}

\section{Research attempt}
\subsection{Literature, source context and scope}
The full original catalogue body and source qualifications are preserved above. The relevant primary workshop passages were read in the cached publisher report. Historical claims and subsequent attributed results are distinguished from the mathematical work below. A restricted example, standard lemma or known theorem is not counted as new progress on the entire original target. Missing assumptions and repeated or rhetorical extracts are not silently turned into different solved problems.

The source combines four questions: cost-based sensor/actuator selection, distributed controllability tests, changes to dynamics or communication links, and PDE-ADMM convergence. They are not one specified optimization problem. This attempt isolates the link-addition question in a finite-dimensional discrete-time model and tests an appealing monotonicity inference.
\subsection{Adding a dynamics link can destroy controllability}
Consider $x_{t+1}=Ax_t+Bu_t$ with
\[
 A_0=\begin{pmatrix}1/4&0\\0&1/2\end{pmatrix},\quad
 A_1=\begin{pmatrix}1/4&1/4\\0&1/2\end{pmatrix},\quad B=\begin{pmatrix}1\\1\end{pmatrix}.
\]
Only one previously zero dynamics entry is added; both matrices have eigenvalues $1/4,1/2$ and are stable.
\begin{proposition}
The first system is controllable in two steps, while the second is not controllable in any number of steps. Their two-step Gramian determinants are $1/16$ and zero, respectively.
\end{proposition}
\begin{proof}
The columns of $[B,A_0B]$ are $(1,1)^T,(1/4,1/2)^T$, so their determinant is $1/4$. Arbitrary target states from zero can be reached with two controls because this matrix (up to column order) is invertible. For $A_1$, one has $A_1B=B/2$, hence $A_1^tB=2^{-t}B$ for every $t\ge0$; the reachable space is only the line spanned by $B$. The two-step Gramian equals $CC^T$ with $C=[B,AB]$, so its determinant is $(\det C)^2$, as asserted.
\end{proof}
\subsection{Where a monotonicity proof actually works}
Keep $A$ fixed and add independent actuator columns, replacing $B$ by $[B,D]$. For any integer horizon $T\ge1$,
\[
 W_T([B,D])-W_T(B)=\sum_{t=0}^{T-1}A^tDD^T(A^T)^t\succeq0,
 \quad W_T(B)=\sum_{t=0}^{T-1}A^tBB^T(A^T)^t.
\]
Indeed each summand has quadratic form $\|D^T(A^T)^tv\|^2\ge0$. If $W_T(B)$ is positive definite, writing the new Gramian as $W_T(B)^{1/2}(I+K)W_T(B)^{1/2}$ with $K\succeq0$ proves that its determinant is at least the old determinant: every eigenvalue of $I+K$ is at least one. If the old Gramian is singular, its determinant is zero and the new one is nonnegative. This conclusion applies to genuinely additional unconstrained inputs; it does not justify changing $A$ or relaxing one constraint while changing disturbances.
\section{Verification and remaining gaps}
Exact rational checks reproduce both controllability matrices, determinants and the added-actuator PSD difference. The failed inference is transferring actuator monotonicity to new dynamics links. Distributed communication, co-design budgets and costs, disturbance propagation and PDE-ADMM convergence remain untreated.

\section{References and source attribution}
Primary workshop report: \url{https://doi.org/10.4171/owr/2024/56}. Relevant passages read: OWR 2024/56, printed pages 3239--3240, Conger and four related open questions.

All retrievals were checked on 11 October 2026. Original question attribution remains in the preserved Sources section. External theorems are attributed and not claimed as new discoveries or fully verified merely by citation.


\paragraph{Integration review.} Codex (GPT-6; exact serving revision not exposed) reviewed the written arguments and expanded layout and source-reference wrappers. Repository metadata and the production web renderer were checked using Python 3.13 and JavaScript with MathJax 3.2.2. This model-assisted review is not independent human proof verification; no PDF compilation or formal proof checking is asserted.

\section{Completion Estimate}
\noindent\textbf{COMPLETION: 0\%}
No substantiated new progress toward the entire original target is claimed. The worked results and failed route are delimited above; a known result, illustrative model, or auxiliary obstruction is not counted as original-target completion.

\section{Final status and attribution}
\textbf{Attempt status: unresolved.} This describes this research attempt, not a new ruling on the source catalogue's status or on external literature claims.
The mathematics and this self-audit were generated by OpenAI \texttt{gpt-6-astra}, reasoning \texttt{ultra}, on 11 October 2026. Codex (GPT-6; exact serving revision not exposed) orchestrated the batch, read and reviewed the written mathematical arguments, independently reran the finite checks, and contributed formatting/repository review. No human mathematical authorship, independent human proof checking, or formal verification is claimed. The historical preparation paragraphs in the reproduced catalogue body concern that body only.

Contribution prompt A was applied to the complete definition and supplied source record, followed by the adversarial checks of prompt B. Local reproduction files are \texttt{batch026/checks.py}, \texttt{computations.json}, and the per-ID audit JSON. They were not public links at preparation time. Numerical checks are finite and do not prove unrestricted results. No LaTeX engine was found locally; compilation was not performed by this research agent.

\end{document}
